\documentclass{erdc}

\usepackage{amsmath}
\usepackage{booktabs}
\usepackage{tabularx}

% The pinned TeX image does not ship ERDC's historical mathgifg font package.
% A local compatibility shim keeps the class's documented interface while
% allowing the report to use the standard Computer Modern math fonts.

\reportnum{NRU-TR-26-014}
\program{Urban Systems and Service Reliability}
\title{A Field Guide to Reliable Public-Service Pilots}
\subtitle{Measurement, handover, and lessons from a six-month trial}
\date{August 2026}
\author{Maya Chen \and Idris Okafor}
\affiliation{Northbridge Research Unit\\
  14 Meridian Square\\
  Northbridge NB4 2QJ}
\reporttype{Technical Information Report}
\additionalinfo{Prepared for the Public Service Reliability Programme}
\contractnum{NRU-26-C-014}
\monitoredby{Northbridge Research Unit\\
  Programme Evaluation Office}

\begin{document}
\frontmatter
% The historical ERDC cover always requests the USACE castle artwork.  Keep
% this generic starter independent of that institutional mark while retaining
% the cover-art slot's proportions.
\let\erdcOriginalIncludegraphics\includegraphics
\renewcommand{\includegraphics}[2][]{\rule{65pt}{40pt}}
\maketitle
\let\includegraphics\erdcOriginalIncludegraphics
\tableofcontents
\listoffiguresandtables

\chapter{Executive Summary}

This report records a six-month pilot for improving the reliability of a public
service that is delivered through several local teams. The design is intentionally
modular: each team can adopt the measurement sheet, weekly review, and handover
checklist without changing its operating system.

The pilot reduced the median time from a reported fault to a confirmed owner from
31 hours to 9 hours. The result depended less on new equipment than on making
ownership visible and closing the loop with the person who raised the issue.

\begin{table}[ht]
  \centering
  \caption{Pilot measures}
  \begin{tabular}{@{}lrrr@{}}
    \toprule
    Measure & Baseline & Pilot & Change \\
    \midrule
    Median ownership time (hours) & 31 & 9 & $-71\%$ \\
    Items closed within target & 54\% & 82\% & $+28$ pp \\
    Repeat reports within 30 days & 19\% & 12\% & $-7$ pp \\
    \bottomrule
  \end{tabular}
\end{table}

\chapter{Pilot Design}
\section{Operating model}

Every report entered a shared queue with three mandatory fields: location,
severity, and named owner. A short morning review handled new high-severity
items; a weekly review looked for repeat failures and work that had crossed its
target date.

\section{Measurement}

For each item $i$, the pilot measured the elapsed ownership time
\begin{equation}
  T_i = t_{\mathrm{owner},i} - t_{\mathrm{report},i}.
\end{equation}
The headline result is the median of $T_i$, not the mean, because a small number
of unusually complex items would otherwise dominate the operational signal.

\chapter{Results}

The largest improvement occurred in the first eight weeks, when the team adopted
a visible queue and assigned a duty manager for each shift. Improvement slowed
after the queue became routine, suggesting that the next phase should focus on
preventing repeat reports rather than reducing ownership time still further.

\begin{table}[ht]
  \centering
  \caption{Weekly pilot summary}
  \begin{tabularx}{0.9\textwidth}{@{}l r r X@{}}
    \toprule
    Week & New items & Closed & Observation \\
    \midrule
    1--2 & 38 & 29 & Establish baseline and duty-manager rotation \\
    3--4 & 41 & 38 & Queue ownership becomes visible \\
    5--8 & 73 & 67 & Repeat faults become the main review topic \\
    9--12 & 64 & 61 & Handover checklist adopted by all teams \\
    \bottomrule
  \end{tabularx}
\end{table}

\chapter{Handover Checklist}

Before the pilot closes, the receiving team should confirm that it has:
\begin{itemize}
  \item an owner and deputy for each service area;
  \item the baseline and target measures, with definitions attached;
  \item an escalation route for overdue and high-severity items; and
  \item a monthly review date for repeat failures and unresolved risks.
\end{itemize}

\chapter{Limitations}

The pilot was not a randomised experiment and the reporting process changed during
the first two weeks. Results should therefore be treated as operational evidence,
not as a causal estimate that can be transferred unchanged to another service.

\chapter*{Document control}
Version 1.0; prepared as an editable technical-report starter. Replace the
identifiers, programme names, measurements, and contacts before publication.
\end{document}
